\section{Most Stones Removed with Same Row or Column}
\label{sec:graph-most-stones-removed}

\problemstatement{Given the coordinates of $n$ stones on a 2D plane, a
stone can be \textbf{removed} if some \textbf{other}, not-yet-removed
stone shares its row or its column. Find the maximum number of stones
that can be removed.}

\begin{patternbox}
\emph{``Stones connected by shared rows/columns''} is a Disjoint Set
grouping problem in disguise: union every pair of stones sharing a row
or column, then within each resulting connected group, \textbf{every
stone but one} can always be removed (remove them one at a time,
always leaving at least one survivor in the group to justify removing
the next). The answer is simply
$\textit{totalStones} - \textit{numberOfGroups}$.
\end{patternbox}

\subsection*{Why one survivor per group is always achievable}

Within a connected group of stones (connected via a chain of shared
rows/columns), pick any spanning tree of that connectivity (exactly
as in Section~\ref{sec:graph-kruskal}'s spanning tree idea) and remove
stones as \textbf{leaves} of that tree, working inward: a leaf stone
always shares its row or column with its tree-parent, which is
guaranteed not yet removed (it's deeper in the removal order), so the
removal condition is always satisfiable until only the root remains.
This is why the answer doesn't depend on \emph{how} stones happen to
connect within a group -- only on \emph{how many} groups there are.

\subsection*{Building the Disjoint Set over rows and columns}

Rather than unioning stones pairwise directly (which would need
checking every pair, $O(n^2)$), union each stone's \textbf{row} and
\textbf{column} together in a combined coordinate space: treat row
indices and column indices as Disjoint Set elements in two disjoint
ranges (e.g., column $c$ mapped to $c + 10001$ to avoid colliding with
row indices). Two stones sharing a row or column end up in the same
Disjoint Set group automatically, transitively, without ever comparing
stones to each other directly.

\subsection*{C++ implementation}

\begin{lstlisting}
#include <bits/stdc++.h>
using namespace std;

class DisjointSet {
    unordered_map<int,int> parent;
public:
    int find(int x) {
        if (parent.find(x) == parent.end()) parent[x] = x;
        if (parent[x] == x) return x;
        return parent[x] = find(parent[x]);
    }
    void unite(int x, int y) {
        parent[find(x)] = find(y);
    }
};

int removeStones(vector<vector<int>>& stones) {
    DisjointSet ds;
    for (auto& s : stones) {
        ds.unite(s[0], s[1] + 10001); // column shifted into a separate range
    }

    unordered_set<int> roots;
    for (auto& s : stones) {
        roots.insert(ds.find(s[0]));
    }
    return (int)stones.size() - (int)roots.size();
}

int main() {
    vector<vector<int>> stones = {{0,0},{0,1},{1,0},{1,2},{2,1},{2,2}};
    cout << removeStones(stones) << endl; // 5
    return 0;
}
\end{lstlisting}

\subsection*{Dry run}

\dryrun{Take $\textit{stones}=[(0,0),(0,1),(1,0),(1,2),(2,1),(2,2)]$.}

Union (row $0$, col $0{+}10001$), (row $0$, col $1{+}10001$), (row
$1$, col $0{+}10001$), (row $1$, col $2{+}10001$), (row $2$, col
$1{+}10001$), (row $2$, col $2{+}10001$). Tracing connectivity: stone
$(0,0)$ links row $0$ and col $10001$; stone $(1,0)$ links row $1$ and
col $10001$ -- so row $0$ and row $1$ end up in the \emph{same} group
via shared column $0$. Continuing, every row and column index here
ends up merged into a \textbf{single} group (rows $0,1,2$ and columns
$0,1,2$ all connected through the chain of shared coordinates). All
$6$ stones' row-roots resolve to this one group: $\textit{roots}$ has
size $1$. Answer: $6 - 1 = 5$.

\begin{complexitybox}
\textbf{Time Complexity.} $O(n \cdot \alpha(n))$ -- unioning and
finding over $n$ stones.
\textbf{Space Complexity.} $O(n)$ for the Disjoint Set's hash-map-based
parent storage (rows and columns are sparse, so a hash map avoids
allocating a dense array over the full coordinate range).
\end{complexitybox}

\begin{takeawaybox}
Mapping two conceptually different categories (rows and columns) into
a single Disjoint Set by shifting one category's indices into a
disjoint numeric range is a general technique whenever ``connected
through a shared attribute'' needs to be modeled without explicitly
enumerating every pairwise connection.
\end{takeawaybox}
